\documentclass[11pt]{article}

\usepackage[a4paper,margin=2.5cm]{geometry}
\usepackage[T1]{fontenc}
\usepackage{lmodern}
\usepackage{xcolor}
\usepackage{fancyhdr}
\usepackage{titlesec}
\usepackage{listings}
\usepackage{booktabs}
\usepackage{array}
\usepackage{tabularx}
\usepackage{enumitem}
\usepackage{hyperref}

\definecolor{ratlemaroon}{RGB}{130,20,30}
\definecolor{ratlegray}{RGB}{90,90,90}
\definecolor{ratlegreen}{RGB}{40,110,60}
\definecolor{ratleframe}{RGB}{190,190,190}

\hypersetup{
  colorlinks=true,
  linkcolor=ratlemaroon,
  citecolor=ratlemaroon,
  urlcolor=ratlemaroon,
  pdftitle={CosmosLog Mission},
  pdfauthor={Carla Gonzalez Mina}
}

\lstdefinelanguage{Rascal}{
  morekeywords={module,import,extend,syntax,lexical,keyword,layout,start,data,public,%
    test,void,str,int,bool,loc,list,set,map,tuple,value,default,visit,case,switch,%
    if,else,for,while,return,true,false},
  sensitive=true,
  morecomment=[l]{//},
  morestring=[b]"
}

\lstdefinelanguage{EBNF}{
  morekeywords={},
  sensitive=true,
  morecomment=[l]{\#}
}

\lstset{
  basicstyle=\ttfamily\footnotesize,
  keywordstyle=\color{ratlemaroon}\bfseries,
  commentstyle=\color{ratlegray}\itshape,
  stringstyle=\color{ratlegreen},
  numbers=left,
  numberstyle=\tiny\color{ratlegray},
  numbersep=8pt,
  frame=single,
  rulecolor=\color{ratleframe},
  framesep=4pt,
  breaklines=true,
  breakatwhitespace=false,
  captionpos=b,
  tabsize=2,
  columns=fullflexible,
  showstringspaces=false,
  xleftmargin=1em,
  xrightmargin=1em
}

\setlength{\headheight}{34pt}
\pagestyle{fancy}
\fancyhf{}
\renewcommand{\headrulewidth}{0.4pt}
\fancyhead[R]{\footnotesize
  ISIS-2111 PLE\\
  Project 1 Review Activity\\
  Date: \today
}
\fancyfoot[C]{\thepage}

\titleformat{\section}{\large\bfseries}{\thesection}{0.6em}{}
\titleformat{\subsection}{\normalsize\bfseries}{\thesubsection}{0.6em}{}

\newcommand{\coverheader}{%
  \begin{minipage}[t]{0.55\textwidth}
    {\Large\bfseries Universidad de los Andes}\\[2pt]
    {\normalsize Department of Systems and Computing Engineering}
  \end{minipage}%
  \hfill
  \begin{minipage}[t]{0.4\textwidth}
    \raggedleft
    ISIS-2111 PLE\\
    Project 1 Review Activity\\
    Date: \today
  \end{minipage}
  \par\noindent\rule{\textwidth}{0.6pt}
}

\title{\vspace{-2.5cm}}
\date{}

\begin{document}

\thispagestyle{empty}
\coverheader

\vspace{1.2cm}
\begin{center}
  {\LARGE\bfseries CosmosLog Mission}\\[6pt]
  {\large Design and Validation of a Review Activity for the\\ Objectilang Grammar Assignment}\\[10pt]
  {\normalsize Prepared by Carla Gonzalez Mina}
\end{center}

\vspace{0.6cm}

\noindent
This report documents the design, the teaching materials, and the technical
validation of a ninety minute review session built for ISIS-2111 Programming
Language Essentials. The session prepares students for Task 1 of Project 1,
writing the EBNF grammar of Objectilang, by having them run the same
procedure end to end on a different, smaller language first.

\vspace{0.4cm}
\tableofcontents

\newpage

\section{Introduction}

Project 1 asks each team to submit a complete EBNF grammar for Objectilang,
an object oriented language with generics, covariance, and contravariance,
described in the course handout through a small set of example programs.
The \emph{Rascal DSL Tutorial} teaches the mechanics needed to eventually
turn a grammar like that into a working parser: concrete syntax rules,
lexical rules, keyword exclusion, and the operators that express repetition,
optionality, and choice.

Both documents assume a level of comfort with grammar design that most
students have not yet practiced under time pressure. The review activity
described here closes that gap. Students spend ninety minutes reconstructing,
repairing, translating, and testing a grammar for a language unrelated to
Objectilang, so that every mistake they are going to make in Project 1 gets
made once already, in a lower stakes setting, with a partner and a teaching
assistant in the room.

\section{Diagnostic Analysis}

\subsection{What Project 1 Requires}

Task 1 asks for the alphabet of terminal symbols, the set of nonterminals,
the distinguished start symbol, and the production rules in EBNF, where
nonterminals are written as bare words and terminals appear in quotes. The
grammar has to account for object declarations, inheritance through
\texttt{extends}, values and variables, methods and method calls, basic
types, arithmetic, boolean, and string expressions, \texttt{if-then-else}
and \texttt{for}, and generic type parameters carrying an optional
covariance or contravariance modifier. Reserved words, identifiers, and the
operator and punctuation alphabet are given directly in the handout.

\subsection{What the Tutorial Contributes}

Section 4 of the tutorial shows how an EBNF-shaped idea becomes a Rascal
\texttt{syntax} rule: \texttt{'terminal'} literals, named nonterminals, and
the \texttt{+}, \texttt{*}, \texttt{?}, and \texttt{|} operators, together
with \texttt{lexical} rules for things like identifiers and integers, and a
\texttt{keyword} declaration that excludes reserved words from the
identifier alphabet. Later sections cover abstract syntax, three styles of
code generation, and validation, but the review activity only needs Section
4, since Project 1 itself asks for a grammar on paper, not a running parser.

\subsection{What the Skeleton Actually Contains}

The folder distributed alongside the assignment, \texttt{Codigo
Ordenado-2}, was inspected file by file before any activity was designed
around it. It does not contain a project skeleton for Objectilang. It holds
the tutorial's own reference code, organized by section, built around the
tutorial's example language of Planning, PersonTasks, and Task. It has no
\texttt{META-INF/RASCAL.MF}, no \texttt{src/main/rascal} layout, and no
content related to Objectilang. That code was still useful: it is a fully
working Rascal project in a domain other than Objectilang, exactly the kind
of material the activity needed for hands on practice without handing
students a solution to their actual assignment.

\subsection{Anticipated Difficulties}

Four difficulties show up reliably when students first design a grammar.
Confusing a terminal with a nonterminal, especially when a keyword is not
quoted. Choosing the wrong repetition or optionality operator, most often
because the operator was inferred from what feels natural rather than from
what the examples actually show. Forgetting that identifiers must exclude
reserved words, an omission that produces a grammar that still accepts
every positive example handed to it, which makes the bug invisible until
someone tries a case built to expose it. Representing a generic type
parameter with an optional variance modifier, the hardest single rule in
the whole assignment, since it requires composing an optional element with
a choice inside a declaration.

\section{Activity Design}

\subsection{Learning Objectives}

By the end of the session a team should be able to tell terminals,
nonterminals, the start symbol, and production rules apart in an unfamiliar
grammar; turn a set of valid example programs into a complete EBNF grammar;
find and justify errors in a grammar written by someone else; translate a
small EBNF rule into Rascal concrete syntax; tell a \texttt{syntax} rule
from a \texttt{lexical} rule and a reserved word declaration; design test
cases, both positive and negative, that would expose a specific class of
bug; and represent a generic type parameter with an optional variance
modifier, without having seen Objectilang's own version of that rule.

\subsection{Format}

Ninety minutes, pairs or small groups, one computer per group already
running a Rascal enabled VS Code, and a teaching assistant circulating
through the room. Groups work through six missions in order, each with a
time limit, up to three hints available on request, and a defined bar for
clearing the mission.

\subsection{Domain Choice}

The main practice domain is CosmosLog, a small language for space mission
logs, invented for this activity. It shares the shape of the tutorial's own
Planning and Task language, one or more top level entries each containing
one or more nested entries, several kinds of leaf record selected by a
keyword, and an optional trailing field, but with different vocabulary and
a different concrete structure, so that reconstructing its grammar is a
genuine exercise rather than a transcription of something already seen.

\subsection{Mission Overview}

\begin{table}[h!]
\centering
\small
\begin{tabularx}{\textwidth}{@{}l l X@{}}
\toprule
\textbf{Mission} & \textbf{Time} & \textbf{What it trains} \\
\midrule
1. Fast classification & 7 min & Telling terminals, nonterminals, and EBNF operators apart, using a grammar the students already built \\
2. Grammar archaeology & 15 min & Writing a complete EBNF grammar from valid example programs \\
3. Repair bay & 13 min & Finding and correcting six planted bugs in someone else's grammar \\
4. Translation to the engine & 15 min & Turning EBNF into Rascal concrete syntax and getting real parser feedback \\
5. Test case detectives & 10 min & Predicting acceptance or rejection, then designing cases that reveal bugs \\
6. The variance challenge & 10 min & Representing an optional covariance or contravariance modifier in EBNF \\
\bottomrule
\end{tabularx}
\caption{The six missions, their time budget, and the skill each one trains.}
\end{table}

An introduction of five minutes opens the session, and a wrap up discussion
of eight minutes followed by a seven minute closing reflection closes it,
bringing the total to ninety minutes exactly.

\section{The CosmosLog Domain-Specific Language}

\subsection{Reference Grammar}

Listing~\ref{lst:ebnf} is the reference EBNF grammar for CosmosLog, the one
a correct Mission 2 submission should approximate. It is never shown to
students before the session; they reconstruct it themselves from the
example log in Listing~\ref{lst:cosmos-example}.

\begin{lstlisting}[language=EBNF,caption={CosmosLog reference grammar, EBNF form},label={lst:ebnf}]
MissionControlLog ::= Mission+

Mission ::= "Mission:" ID Nave+

Nave ::= "Ship" ID "crew" INT Report+

Report ::= "Report" ReportKind "severity" ":" INT Note?

ReportKind ::= FuelReport | AnomalyReport | CommReport | DockReport

FuelReport ::= "Fuel" INT "percent"
AnomalyReport ::= "Anomaly" STRING
CommReport ::= "Comm" ID
DockReport ::= "Dock" ID

Note ::= "note" ":" STRING
\end{lstlisting}

\subsection{Rascal Concrete Syntax}

Listing~\ref{lst:rascal-grammar} is the same grammar translated into Rascal,
matching the pattern the tutorial teaches in Section 4.1. This is the
grammar Mission 4 asks students to arrive at on their own, starting from a
file that gives them only the lexical alphabet and a set of labeled gaps.

\begin{lstlisting}[language=Rascal,caption={CosmosLog grammar in Rascal concrete syntax},label={lst:rascal-grammar}]
start syntax MissionControlLog
    = missionControlLog: Mission+ missions
;

syntax Mission
    = mission: 'Mission:' ID name Nave+ naves
;

syntax Nave
    = nave: 'Ship' ID name 'crew' INT crewSize Report+ reports
;

syntax Report
    = report: 'Report' ReportKind kind 'severity' ':' INT sev Note? note
;

syntax ReportKind
    = fuel: FuelReport fuelReport
    | anomaly: AnomalyReport anomalyReport
    | comm: CommReport commReport
    | dock: DockReport dockReport
;

syntax FuelReport = fuelReport: 'Fuel' INT pct 'percent' ;
syntax AnomalyReport = anomalyReport: 'Anomaly' STRING description ;
syntax CommReport = commReport: 'Comm' ID target ;
syntax DockReport = dockReport: 'Dock' ID station ;

syntax Note = note: 'note' ':' STRING text ;

keyword Reserved = "Mission" | "Ship" | "crew" | "Report" | "severity"
                  | "note" | "Fuel" | "percent" | "Anomaly" | "Comm"
                  | "Dock" | ":";
\end{lstlisting}

\subsection{Example Instance}

\begin{lstlisting}[caption={A valid CosmosLog log, used throughout Missions 2, 4, and 5},label={lst:cosmos-example}]
Mission: Artemis9
Ship Falcon crew 6
Report Fuel 82 percent severity: 2
Report Comm GroundControl severity: 3

Mission: Odyssey2
Ship Odyssey crew 3
Report Dock StationAlpha severity: 1
Report Fuel 45 percent severity: 5 note: "Refuel scheduled next orbit"

Mission: Voyager3
Ship Nomad crew 12
Report Anomaly "Sensor drift on panel B" severity: 7 note: "Monitoring closely"
Ship Comet crew 2
Report Fuel 10 percent severity: 9 note: "Critical - refuel immediately"
\end{lstlisting}

This single log exercises every construct in the grammar: more than one
mission per file, more than one ship per mission, all four report kinds,
and a report with and without its optional note.

\section{Mission 6 and the Objectilang Connection}

\label{sec:mission6}
Mission 6 is the deliberate bridge to the hardest part of Project 1. It
asks students to design an EBNF rule for \texttt{Manifest}, a generic cargo
container carrying a single type parameter with an optional covariance or
contravariance modifier, in the same spirit as \texttt{List[+A]} and
\texttt{Transformer[-In,+Out]} from the Project 1 handout, without ever
showing them Objectilang's own grammar. One valid answer is:

\begin{lstlisting}[language=EBNF,caption={A valid Mission 6 answer for the Manifest rule}]
Manifest ::= "Manifest" "[" TypeParam "]"
TypeParam ::= Variance? ID
Variance ::= "+" | "-"
\end{lstlisting}

Students are then asked, without writing the rule out, to explain how this
would extend to two type parameters with independent modifiers, which is
exactly the shape of \texttt{Transformer[-In,+Out]}. A grading rubric for
this mission accepts any structurally sound answer, not only the one shown
above, since the point is transferring the pattern, not matching a fixed
string.

\section{Materials Produced}

The activity ships as a self contained folder, \texttt{cosmoslog-activity},
split into a student facing half and an instructor only half.

\begin{itemize}[nosep]
  \item \texttt{FOR\_STUDENTS/student\_handout.md}, the complete session
    handout: context, rules, all six missions, the example log, and a
    submission checklist.
  \item \texttt{FOR\_STUDENTS/cosmoslog-dsl/}, a working Rascal project
    with the layout and lexical alphabet already filled in and the seven
    syntax rules left as labeled gaps for Mission 4, plus the example and
    test instance files used in Missions 2 and 5.
  \item \texttt{FOR\_TA/session\_script.md}, a minute by minute script
    covering timing, instructions, files used, and what to watch for in
    each mission.
  \item \texttt{FOR\_TA/ta\_guide.md}, the answer key: solutions, three
    tiered hints per mission, guiding questions, common mistakes, and a
    short formative rubric.
  \item \texttt{FOR\_TA/reference\_solutions/}, the finished grammar, the
    Mission 2 reference grammar, the Mission 3 bug list with evidence, the
    Mission 6 sample answers, and a small automated regression project
    described in Section~\ref{sec:verification}.
\end{itemize}

\section{Technical Integration and Verification}
\label{sec:verification}

Every technical claim in this report was checked against a running Rascal
installation rather than left as an assumption.

\subsection{Compilation Check}

The reference grammar was compiled with the real Rascal type checker,
invoked through the \texttt{rascal-maven-plugin} \texttt{compile} goal.
The build reported zero errors and zero undefined symbols, confirming that
the grammar is valid Rascal and that \texttt{start[MissionControlLog]}
resolves correctly.

\subsection{Automated Regression Test}

A small self contained project, \texttt{FOR\_TA/reference\_solutions/%
smoke-test/}, bundles the reference grammar together with the example log
and the six Mission 5 test cases, and a \texttt{SmokeTest} module that
parses each one and compares the result against its expected outcome. This
is not a one time check; it can be rerun after any change to the reference
grammar, or against a different Rascal distribution, in under two minutes.

\begin{table}[h!]
\centering
\small
\begin{tabular}{@{}l l l@{}}
\toprule
\textbf{Case} & \textbf{Expected} & \textbf{What it tests} \\
\midrule
Positive log & Accepted & The full example log, all four report kinds \\
E1 & Accepted & A minimal single mission, single ship instance \\
E2 & Accepted & Two missions, a Dock report, no note anywhere \\
E3 & Rejected & Missing the leading \texttt{"Mission:"} literal \\
E4 & Rejected & A string given where \texttt{severity} expects an integer \\
E5 & Rejected & A reserved word used as a ship name \\
E6 & Accepted & A negative crew size, syntactically legal, semantically absurd \\
\bottomrule
\end{tabular}
\caption{The seven automated test cases and their confirmed outcome.}
\end{table}

Every one of the seven cases matched its expected outcome on the last run,
across two independent execution paths: the Rascal interpreter invoked
directly, and the type checker invoked through Maven. Case E6 is kept
deliberately even though it passes, since it is the anchor for the closing
discussion about the difference between a syntax error and a validation
error, the subject of Section 8 of the tutorial.


\section{Mapping Missions to Project 1}

\hspace{1cm}

\begin{table}[h!]
\centering
\small
\begin{tabularx}{\textwidth}{@{}l X@{}}
\toprule
\textbf{Mission} & \textbf{Where it resurfaces in Project 1} \\
\midrule
1 & The requirement to state the sets of terminals and nonterminals explicitly \\
2 & Turning Snippets 1 through 4 into the Objectilang grammar \\
3 & Self review before submitting, catching the same classes of bug \\
4 & Preparation for any future assignment that asks for a working implementation \\
5 & Designing test cases against one's own grammar \\
6 & \texttt{List[+A]} and \texttt{Transformer[-In,+Out]}, see Section~\ref{sec:mission6} \\
\bottomrule
\end{tabularx}
\caption{How each mission maps back onto a concrete requirement of Project 1.}
\end{table}


\section{Conclusion}

The activity turns four abstract skills, reading examples into a grammar,
catching grammar bugs, translating EBNF into a working parser, and
designing test cases, into something a pair of students can practice once,
under time pressure, before those same skills matter for a grade. Every
grammar, every test case, and every expected result in this report was
verified against a running Rascal installation rather than asserted, and
the automated regression project in
\texttt{FOR\_TA/reference\_solutions/smoke-test/} lets that verification be
repeated at any time.

\end{document}
